\documentclass[10pt,a4paper]{amsart}
\usepackage[margin=19mm]{geometry}
\usepackage{amsmath,amssymb,mathtools}
\usepackage[scr=rsfs,cal=euler]{mathalfa}
\usepackage{xfrac}
\usepackage[unicode,hidelinks,bookmarks=false]{hyperref}
\newcommand{\ie}{i.e.\ }
\newcommand{\cf}{cf.\ }
\newcommand{\resp}{respectively\ }
\newcommand{\Subsection}{\subsection}
\providecommand{\nobreakdash}{\nobreak\hbox{-}}
\setlength{\emergencystretch}{2em}
\multlinegap0pt
\usepackage{bm}
\usepackage{bbm}
\usepackage{dsfont}
\usepackage{xspace}
\usepackage{cancel}
\usepackage{mathtools}
\usepackage{amsthm}
\makeatletter
\@ifundefined{thm}{\newtheorem{thm}{Thm}}{}
\@ifundefined{theorem}{\newtheorem{theorem}[thm]{Theorem}}{}
\makeatother
\title[On some open problems in commutative algebra resolved by Ret]{On some open problems in commutative algebra resolved by Rethlas}
\author{Jiedong Jiang, Yixiao Li, Zeming Sun, Yuefeng Wang, Liang Xiao, Jiahong Yu}
\date{Source: arXiv:2605.25259v2 (CC BY 4.0)}
\begin{document}
\maketitle

\noindent\emph{Source and licence.} On some open problems in commutative algebra resolved by Rethlas. Version arXiv:2605.25259v2 (CC BY 4.0), distributed under the Creative Commons Attribution 4.0 International licence, as recorded in the arXiv licence metadata for this exact version and shown on the abstract page https://arxiv.org/abs/2605.25259v2. Full rights notice: licenses/arxiv-2605.25259-CC-BY-4.0.txt.\par\smallskip

\noindent\textit{Editorial scope.} Counted unit: the Theorem whose source label is \texttt{thm:bs62} in the article (source0.tex lines 1084--1091, source label \texttt{thm:bs62}), delivered together with the article's own complete original proof of it (source0.tex lines 1093--1181, one \texttt{proof} environment of 89 lines). No mathematics is written, repaired or re-derived: statement and proof are the article's own text line for line. Required context retained uncounted: none. Every definition, display and label of those lines is retained unchanged. Omitted from this package and not redistributed: 1--1083, 1092--1092, 1182--1278. Nothing inside a retained excerpt is omitted or altered: every retained body is delivered exactly as the article's lines are written, so no omission receipt of any kind accompanies this package. Citations: the retained excerpts carry no citation command at all, so no bibliography entry of the article is transcribed and no reference is redistributed. Reference adaptation: none was needed and none was made; every reference inside the delivered unit resolves inside it, so no unresolved reference or citation is produced. Printed slips kept literal: no printed word, symbol, hypothesis or number of the counted statement or of its proof was altered. Layout and class: the article's own class is \texttt{amsart} with packages including geometry, amssymb, amsmath, mathtools, mathrsfs, enumerate, xcolor, hyperref, cleveref; the wrapper is the lane's pinned \texttt{amsart} editorial wrapper at 10pt on A4, which restates the theorem-like environments the retained text needs and adds the article's own macro definitions that the retained text needs (none), with extra wrapper packages (bm, bbm, dsfont, xspace, cancel, mathtools). The delivered numbering is the unit's own. Assets: no retained span contains a figure environment, a graphics-inclusion command, a PDF-page inclusion, a table or a TikZ picture; no image file is copied into this package, the package declares no asset dependency and no image file is redistributed with it. Licence: version arXiv:2605.25259v2 of this article is distributed under the Creative Commons Attribution 4.0 International licence, as recorded in the arXiv licence metadata for this exact version and shown on the abstract page https://arxiv.org/abs/2605.25259v2. A full rights notice is delivered at licenses/arxiv-2605.25259-CC-BY-4.0.txt. Characters: every retained excerpt is pure ASCII, so the delivered engine, XeLaTeX with its default Latin Modern text font, reports no missing character.
\begin{theorem}\label{thm:bs62}
Let $k$ be a field. There exist an integer $n\ge 1$, a degree sequence
$(d_0,\dots,d_n)$, and an integral point on the corresponding pure ray
in the Boij--S\"oderberg cone for which no $n$-dimensional
$\mathbb Z_{>0}$-graded Lie algebra $\mathfrak g$ generated in degree~$1$
over $k$ admits a finite length graded module $M$ over $U(\mathfrak g)$
whose Betti table is that integral point.
\end{theorem}

\begin{proof}
We give a counterexample with $n=4$ and
\[
   d=(0,1,4,5,6).
\]

\smallskip
\noindent\emph{Step 1: the primitive integral point.} Let
$\beta=(\beta_0,\beta_1,\beta_2,\beta_3,\beta_4)$ be the coefficients of
a pure table on this ray, so the only nonzero Betti entries are
$\beta_{i,d_i}=\beta_i$. The Herzog--K\"uhl equations for a codimension-$4$
pure resolution read
\[
   \sum_{i=0}^4 (-1)^i \beta_i d_i^m=0,\qquad m=0,1,2,3.
\]
Normalizing $\beta_0=1$, the four equations become
\[
   1-\beta_1+\beta_2-\beta_3+\beta_4=0,\qquad
   -\beta_1+4\beta_2-5\beta_3+6\beta_4=0,
\]
\[
   -\beta_1+16\beta_2-25\beta_3+36\beta_4=0,\qquad
   -\beta_1+64\beta_2-125\beta_3+216\beta_4=0.
\]
Subtracting consecutive equations,
\[
   6\beta_2-10\beta_3+15\beta_4=0,\qquad
   12\beta_2-25\beta_3+45\beta_4=0;
\]
subtracting twice the first from the second gives $\beta_3=3\beta_4$, and
back-substitution yields $2\beta_2=5\beta_4$ and $\beta_1=\beta_4$.
Plugging into the first equation gives $\beta_4=2$, so
\[
   (\beta_0,\beta_1,\beta_2,\beta_3,\beta_4)=(1,2,5,6,2),
\]
which is already integral, hence the primitive integral point on the
ray.

\smallskip
\noindent\emph{Step 2: PBW Hilbert series and divisibility.} Let
$\mathfrak g=\bigoplus_{i\ge 1}\mathfrak g_i$ be a finite-dimensional
positively graded Lie algebra over $k$ with $h_i:=\dim_k\mathfrak g_i$.
Choosing a homogeneous basis and applying PBW, the ordered monomials in
basis vectors form a graded $k$-basis of $U(\mathfrak g)$, so
\[
   H_{U(\mathfrak g)}(t)=\prod_{i\ge 1}(1-t^i)^{-h_i}.
\]
If a finite length graded $U(\mathfrak g)$-module $M$ admits a pure
resolution with Betti numbers $\beta_i$ in degrees $d_i$, then taking
alternating sums of Hilbert series in a graded free resolution gives
$H_M(t)=p_M(t)\,H_{U(\mathfrak g)}(t)$ with $p_M(t)=\sum_i(-1)^i\beta_i
t^{d_i}$. Since $H_M(t)$ is a polynomial,
\[
   p_M(t) \text{ is divisible by } \prod_{i\ge 1}(1-t^i)^{h_i}.
\]

\smallskip
\noindent\emph{Step 3: contradiction.} Suppose $\mathfrak g$ is a
$4$-dimensional positively graded Lie algebra over $k$ generated in
degree~$1$ and $M$ is a finite length graded $U(\mathfrak g)$-module with
the Betti table from Step~1. The alternating Betti polynomial is
\[
   p(t)=1-2t+5t^4-6t^5+2t^6=(1-t)^4\bigl(2t^2+2t+1\bigr).
\]
By Step~2 it is divisible by $\prod_{i\ge 1}(1-t^i)^{h_i}$. Since
$\dim\mathfrak g=\sum h_i=4$, dividing by $(1-t)^4$ leaves
\[
   q(t):=2t^2+2t+1
\]
divisible by $\prod_{i\ge 2}(1+t+\cdots+t^{i-1})^{h_i}$. But:
$q(-1)=1\ne 0$, so $1+t$ does not divide $q$; modulo $1+t+t^2$ one has
$t^2+t\equiv -1$, hence $q(t)\equiv -1$, so $1+t+t^2$ does not divide
$q$ over any field; and $\deg q=2$, so no $1+t+\cdots+t^{i-1}$ with
$i\ge 4$ can divide $q$. Therefore $h_i=0$ for all $i\ge 2$, i.e.
$\mathfrak g$ is concentrated in degree~$1$. Then
$[\mathfrak g,\mathfrak g]\subseteq\mathfrak g_2=0$, so $\mathfrak g$ is
abelian and
\[
   U(\mathfrak g)\cong k[x_1,x_2,x_3,x_4]=:R.
\]
So $M$ would be a finite length graded $R$-module with Betti table
$(1,2,5,6,2)$. But $\beta_0=1$ means $M\cong R/I$ for some homogeneous
ideal $I$, and $\beta_1=2$ means $I$ is minimally generated by two
homogeneous elements; by Krull's height theorem $\operatorname{ht}(I)\le 2$,
so $\dim(R/I)\ge 4-2=2$, contradicting finite length.

Thus no such pair $(\mathfrak g,M)$ exists for the degree sequence
$(0,1,4,5,6)$ and the integral point $(1,2,5,6,2)$.
\end{proof}

\end{document}
